\subsection{N-gram-malli optimointiongelmana}

Luvussa \ref{sec:n-gram} \(n\)-gram-malli esitettiin lukumääriin ja niiden
normalisointiin perustuvana frekvenssimallina. Sama malli voidaan kuitenkin
tulkita myös suurimman uskottavuuden estimointina, mikä tekee yhteyden
myöhempiin parametrisiin kielimalleihin näkyväksi.

Olkoon mallin parametreina jokaiselle kontekstille
\(c = (w_{t-n+1}, \ldots, w_{t-1})\) määritelty todennäköisyysjakauma
\[
  \theta_c = (\theta_{c,1}, \theta_{c,2}, \ldots, \theta_{c,|V|}),
  \qquad
  \sum_{j=1}^{|V|} \theta_{c,j} = 1,
  \qquad
  \theta_{c,j} \ge 0,
\]
missä \(\theta_{c,j} = P_\theta(w_t = v_j \mid c)\). Markov-oletuksen nojalla havaintoaineiston
\(W = (w_1, \ldots, w_{|W|})\) uskottavuus voidaan kirjoittaa muodossa
\[
  P(W;\theta)
  \approx
  \prod_{t=1}^{|W|}
  P_\theta(w_t \mid w_{t-n+1}, \ldots, w_{t-1}).
\]

Kun aineisto ryhmitellään kontekstin \(c\) ja sitä seuraavan tokenin \(v_j\)
mukaan, saadaan
\[
  P(W;\theta)
  =
  \prod_c \prod_{j=1}^{|V|}
  \theta_{c,j}^{C_{c,j}},
\]
missä \(C_{c,j}\) on niiden havaintojen lukumäärä, joissa kontekstia \(c\)
seuraa tokeni \(v_j\).

Koneoppimisessa optimointi esitetään tavallisesti minimointitehtävänä, joten
määritellään häviöksi negatiivinen log-uskottavuus
\[
  J(\theta;W)
  =
  -\log P(W;\theta)
  =
  -\sum_c \sum_{j=1}^{|V|}
  C_{c,j} \log \theta_{c,j}.
\]
Koska parametrit \(\theta_c\) esiintyvät vain omaa
kontekstiaan vastaavissa termeissä, optimointiongelma hajoaa erillisiksi
kontekstikohtaisiksi osatehtäviksi:
\[
  \min_{\theta_c}
  -\sum_{j=1}^{|V|} C_{c,j} \log \theta_{c,j}
  \qquad
  \text{s.e.}
  \qquad
  \sum_{j=1}^{|V|}\theta_{c,j} = 1.
\]

Tämä voidaan ratkaista Lagrangen menetelmällä. Määritellään
\[
  \mathcal{L}_{\mathrm{Lag}}(\theta_c,\lambda)
  =
  \sum_{j=1}^{|V|}
  C_{c,j}\log \theta_{c,j}
  +
  \lambda\left(1-\sum_{j=1}^{|V|}\theta_{c,j}\right).
\]
Derivoimalla \(\theta_{c,j}\):n suhteen saadaan
\[
  \frac{\partial \mathcal{L}_{\mathrm{Lag}}}{\partial \theta_{c,j}}
  =
  \frac{C_{c,j}}{\theta_{c,j}} - \lambda = 0,
\]
mistä seuraa
\[
  \theta_{c,j} = \frac{C_{c,j}}{\lambda}.
\]
Kun tämä sijoitetaan normalisaatioehtoon
\[
  \sum_{j=1}^{|V|}\theta_{c,j} = 1,
\]
saadaan
\[
  \lambda = \sum_{k=1}^{|V|} C_{c,k},
\]
ja siten
\[
  \theta_{c,j}
  =
  \frac{C_{c,j}}{\sum_{k=1}^{|V|} C_{c,k}}.
\]

Tulos osoittaa, että suurimman uskottavuuden estimaatti on täsmälleen sama kuin
luvussa \ref{sec:n-gram} määritelty rivikohtaisesti normalisoitu
frekvenssimatriisi \(F\). Toisin sanoen klassinen frekvenssipohjainen
\(n\)-gram-malli voidaan tulkita parametrisen mallin
maksimiuskottavuusratkaisuna.

Tämä näkökulma tekee samalla näkyväksi mallin keskeisen rajoitteen: jokaiselle
mahdolliselle kontekstille tarvitaan oma erillinen todennäköisyysjakaumansa.
Kuten luvussa \ref{sec:n-gram} nähtiin, kontekstin pituuden kasvaessa parametrien määrä kasvaa nopeasti, jolloin malli
muuttuu harvaksi ja yleistää heikosti harvinaisiin tai havaitsemattomiin
konteksteihin. Tästä syystä on luontevaa siirtyä malleihin, joissa ehdolliset
todennäköisyydet opitaan jaetun parametrisaation avulla jatkuvassa avaruudessa
sen sijaan, että ne talletettaisiin erillisinä taulukoituina jakaumina.

\subsection{Monikerroksinen perseptroni seuraavan tokenin ennustajana}

Seuraavan tokenin ehdollista jakaumaa voidaan
approksimoida parametrisoidulla neuroverkolla. Yksi yksinkertainen ja
historiallisesti keskeinen ratkaisu tähän on monikerroksinen perseptroni
(engl.\ \emph{multilayer perceptron}, MLP).

MLP on eteenpäinsyöttävä keinotekoinen neuroverkko, joka koostuu
syötekerroksesta, yhdestä tai useammasta piilokerroksesta sekä
ulostulokerroksesta. Sen historialliset juuret ovat Frank Rosenblattin
perseptronitutkimuksessa \cite{rosenblatt1958perceptron}\cite{rosenblatt1962principles}. Näissä töissä esitettiin varhaisia
perseptronimalleja ja kerroksellisia rakenteita, mutta tehokas yleiskäyttöinen
menetelmä kaikkien kerrosten parametrien oppimiseen puuttui vielä.

Neuroverkkojen varhaista kehitystä hidasti merkittävästi Marvin Minskyn ja
Seymour Papertin teos, jossa analysoitiin yksikerroksisten
perseptronien rajoitteita\\\cite{minsky1969perceptrons}. Teos osoitti, että yksinkertainen perseptroni ei voi
ratkaista tiettyjä lineaarisesti erottamattomia ongelmia, kuten XOR-tehtävää.
Vaikka kritiikki kohdistui ensisijaisesti yksikerroksisiin malleihin, sen
vaikutus ulottui laajemmin neuroverkkotutkimukseen ja osaltaan vähensi
kiinnostusta monikerroksisiin verkkoihin useiden vuosien ajaksi.

Monikerroksisten verkkojen tehokkaan opettamisen kannalta keskeinen teoreettinen
lähtökohta syntyi Seppo Linnainmaan työssä, jossa esitettiin käänteisen suunnan automaattinen derivointi menetelmä\cite{linnainmaa1970taylor}.
Paul Werbos puolestaan sovelsi tätä ajatusta neuroverkkoihin väitöskirjassaan
ja myöhemmin artikkelissaan, jossa termi vastavirta-algoritmi (engl. \emph{backpropagation}) vakiintui \cite{werbos1974beyond}\cite{werbos1990bptt}.

Monikerroksisten eteenpäinsyöttävien verkkojen läpimurto
tapahtui kuitenkin vasta, kun David Rumelhart, Geoffrey Hinton ja Ronald
Williamsin teos popularisoi vastavirta-algoritmin\cite{rumelhart1986learning}. Tämä työ teki MLP:stä käytännöllisen gradienttipohjaisesti opetettavan funktionapproksimaattorin ja loi perustan
myöhemmälle syväoppimiselle.

Matemaattisesti MLP voidaan ymmärtää funktion approksimaattorina, joka muodostaa
kuvauksen syöteavaruudesta tulosteavaruuteen. Tämä kuvaus toteutetaan
yhdistettynä funktiona, joka koostuu peräkkäisistä affiineista muunnoksista ja
niitä seuraavista epälineaarisista aktivaatiofunktioista. Tämän ansiosta MLP
kykenee mallintamaan sellaisia syötteen ja tulosteen välisiä riippuvuuksia,
joita lineaarinen malli ei voi esittää. Mallin toimintaa ohjaavat opittavat
parametrit, tyypillisesti painomatriisit ja bias-vektorit, joiden arvot
estimoidaan datasta minimoimalla valittua kohdefunktiota numeerisilla
optimointimenetelmillä, kuten gradienttimenettelyllä
\cite{goodfellow2016deep}.

Olkoon \(c\) tarkasteltava konteksti ja
\(x(c)\in\mathbb{R}^{d_0}\) sitä vastaava syötevektori.
Monikerroksinen perceptroni muuntaa tämän syötteen vaiheittain
piirre-esitykseksi soveltamalla peräkkäisiä affiinisia muunnoksia ja
epälineaarisia aktivaatioita. Jos mallissa on \(L\) piilokerrosta,
voidaan sen laskenta esittää muodossa
\begin{align*}
  h^{(0)} &= x(c),\\
  a^{(\ell)} &= W^{(\ell)} h^{(\ell-1)} + b^{(\ell)},
  & \ell = 1,\dots,L, \\
  h^{(\ell)} &= \phi^{(\ell)}\!\left(a^{(\ell)}\right),
  & \ell = 1,\dots,L,
\end{align*}
missä \(W^{(\ell)}\in\mathbb{R}^{d_\ell\times d_{\ell-1}}\) on kerroksen
\(\ell\) painomatriisi, \(b^{(\ell)}\in\mathbb{R}^{d_\ell}\) sen bias-vektori,
\(a^{(\ell)}\in\mathbb{R}^{d_\ell}\) affiinisen muunnoksen tulos ja
\(h^{(\ell)}\in\mathbb{R}^{d_\ell}\) kerroksen \(\ell\) piiloesitys.
Funktio \(\phi^{(\ell)}\) on kerroksen epälineaarinen aktivaatiofunktio.

Viimeisestä piilokerroksesta muodostetaan ulostulon
normalisoimattomat pistemäärät eli logitit
\[
  z = W^{(\mathrm{out})} h^{(L)} + b^{(\mathrm{out})},
\]
missä \(W^{(\mathrm{out})}\in\mathbb{R}^{|V|\times d_L}\),
\(b^{(\mathrm{out})}\in\mathbb{R}^{|V|}\) ja \(z\in\mathbb{R}^{|V|}\).
Vektorin \(z\) komponentti \(z_j\) vastaa sanaston tokenin \(v_j\)
normalisoimatonta pistemäärää. Näistä muodostetaan ehdollinen
todennäköisyysjakauma softmax-funktiolla:
\[
  p_\theta(w_t = v_j \mid c)
  =
  \frac{\exp(z_j)}{\sum_{k=1}^{|V|}\exp(z_k)},
  \qquad j=1,\dots,|V|.
\]

MLP muuntaa kontekstin ensin yhden tai usean kerroksen kautta
abstraktiksi piirre-esitykseksi, jonka perusteella se laskee jokaiselle
sanaston tokenille pistemäärän ja normalisoi nämä pistemäärät
todennäköisyysjakaumaksi.

Mallin oppiminen perustuu samaan negatiiviseen log-uskottavuuden minimointiin kuin \(n\)-gram-mallissa:
\[
  J(\theta)
  =
  -\frac{1}{|W|}
  \sum_{t=1}^{|W|}
  \log p_\theta(w_t \mid w_{t-n+1}, \ldots, w_{t-1}).
\]

Erona perinteiseen frekvenssipohjaiseen \(n\)-gram-malliin, joka estimoi
jokaiselle mahdolliselle diskreetille kontekstille oman erillisen
parametrivektorinsa, MLP approksimoi todennäköisyyksiä jatkuvassa avaruudessa.
Koska erilliset tokenit projisoidaan matalaulotteisiksi tiheiksi
upotusvektoreiksi \(v \in \mathbb{R}^d\), malli oppii esittämään sanojen
semanttisia ja syntaktisia samankaltaisuuksia vektoriesitysten avulla.
MLP:n approksimoima funktio on jatkuva ja tyypillisesti sileä siinä mielessä,
että pienet muutokset syöteavaruudessa voivat johtaa samankaltaisiin
ulostulojakaumiin. Tästä syystä malli kykenee yleistämään myös harvinaisiin tai
aiemmin havaitsemattomiin konteksteihin hyödyntämällä sellaisten
koulutusdatassa esiintyneiden kontekstien rakennetta, joiden esitykset
sijaitsevat lähellä niitä yhteisessä avaruudessa \cite{bengio2003neural}. Tämä
lieventää klassisia \(n\)-gram-malleja vaivaavaa nollatodennäköisyyksien
ongelmaa ilman erillisiä tasoitusmenetelmiä.

Frekvenssipohjaista mallia koskevat tekniset rajoitteet, kuten
eksponentiaalinen muistitarpeen kasvu ja lukumatriisin harvuus, voidaan siten
osittain ratkaista MLP-mallin avulla.
MLP-rakenteessa opittavat painomatriisit ja bias-vektorit jakautuvat kaikkien
kontekstien kesken. Siten jokainen harjoitusnäyte päivittää koko neuroverkon
parametreja, mikä parantaa myös sellaisten lauseiden ennustamista, joita ei ole
esiintynyt opetusdatassa, mutta jotka jakavat samankaltaisia piirteitä
opittujen esitysten kautta. Tämä tekee MLP-pohjaisesta \(n\)-gram-mallista
huomattavasti joustavamman, muistikapasiteetiltaan tehokkaamman ja
yleistyskykyisemmän kuin perinteinen frekvenssiperusteinen
taulukkolukuesitys \cite{bengio2003neural,goodfellow2016deep,jurafsky2023speech}.
